% SCRATCH DRAFT — Vol 2 courtroom coda, Ch2 "Force, Inertia, and Momentum" (task vol2-coda, author: Podo).
% THIS IS NOT THE BOOK. Do not \input. Floor commit not yet banked; 02.tex untouched.
%
% CONTINUITY: state carries from Ch1 (locked). At Ch1's close the record held: one wheel going
% around, one citation entered, the hall-speedometer reading the crank, the charge/trace throughline
% sounded once, the radar gun + GPS named as still-coming. Ch2 opens from THAT — the wheel does
% something new (it speeds and slows) — and does NOT reintroduce the wheel cold.
%
% PROPOSED SEAM in books/experimentation/latex/chapters/02.tex:
%   KEEP locked body through line 262 ("...the residue a closed loop carries.")
%   KEEP the \rule separator block (lines 264-268).
%   REPLACE the OLD beat now at lines 270-275
%     ("A charge, once entered, must be accounted for ... A charge that balances is a charge still
%      standing when it is called.")
%   WITH the prose below. The new arc GROWS the old beat's best sentence
%   ("A charge, once entered, must be accounted for") into the continuing story.
%
% >>>>>>>>>>>>>>>>>>>>>> APPEND BEGINS (after the \rule) >>>>>>>>>>>>>>>>>>>>>>

\noindent A charge, once entered, must be accounted for. The citation opened on a single mark; but no case
is made of one mark left to sit. Once written, the charge is a debt the record has to answer for --- a tally
that cannot quietly go missing between the moment it is set down and the moment it is called. The chapter has
just watched the instrument learn how the record answers for a charge when the motion behind it will not hold
still.

For the wheel is no longer turning at the steady rate we first met it at. It speeds and it slows, and the
needle moves with it. The hall sensor still trips once per pass, but the passes no longer arrive on a fixed
cadence: each comes a little sooner or a little later than the last, and to keep the record consistent the
instrument must re-file its established run of counts onto the new rate, bringing the old order and the new into
agreement at their seam. That re-filing is the instrument's own act --- ours, the reconciliation a changing
motion forces. What it costs is not ours. How stubbornly a given record resists being redirected, how much
re-filing each unit of change exacts, is a standing price the world sets and the instrument can only pay; and
that price, read off the ledger and charged to no one but the world, is the mass. Give two records the same
push and the light one yields in a few reconciliations while the heavy one must be paid out over many before its
old order comes onto the new rate; the instrument has weighed nothing, only counted how dearly each resists the
re-filing. The pushing is the hand's; the rate the world charges for it is not.

Watch the speedometer under this changing motion and it is the needle's swing that carries the new reading ---
the rate of the count itself changing, a bend in the tally read before any force is named. The straight
continuation of the counts is the standard the hand lays down; the gap the motion opens against it is the
world's, curvature read off the record with no force yet spoken. To read a changing
rate the instrument leans on a calibration it did not need while the wheel ran steady: its own time-base, the
clock by which it stamps each pass as it arrives. The tire's circumference and the gear ratio told it how far
one turn carries; the time-base tells it how the rate of turning is changing --- and that clock, like the grid
before it, was laid by a hand before the wheel ever quickened. (What fixes the time-base itself, far beneath
the gear train, in something no clock we wind can set, runs deeper than mechanics reaches; the book keeps that
floor sealed, with the sensor's own scale, until it has the physics to open both.)

Here the record's charge comes fully into its double meaning. Ch1 sounded it once --- the count is a charge,
and the payment, when it comes, is only that same retained turning read out as owed. Now the charge shows the
other half of why it earns the name: it \emph{balances}. Draw a closed line around a stretch of the record and
let the motion cross it, some of the count entering, some leaving. Nothing passes the account uncounted; what
crosses in one side leaves the other, and the books are made to close. The term the instrument carries to keep
that promise keepable is momentum --- position's conjugate partner, the entry whose inflow and outflow it can
balance against the region's change without once opening the region up. The books are the clerk's; that they
are forced to close is the world's. A charge that balances is a charge still standing when it is called. And the trace now spans that boundary: what the record retains of the turning is
not spilled in the re-filing but carried whole across the reconciliation --- the accumulated rotation handed
intact from the old rate to the new, retained entire. The charge the case will answer for is exactly that
conserved trace, the tally the world will not let vanish between its entry and its answer.

A tally kept this well raises the demand that has hung over the record since the first count: one instrument's
word is not yet the record. A reading that changes is harder to force into agreement than one that holds
steady, and the standard rises to meet it. The car's own speedometer, swinging with the motion, will still have
to be met by the two instruments named as coming and made to land on the same moving number --- one changing
charge, reconciled in three independent hands, not merely tallied in one. Their physics waits for the parts
that carry it; the demand is only sharpened here, that a charge which moves must still be made to agree.

One cost in this chapter, though, the instrument records and cannot resolve: the wall where a body holds fast
until a threshold is crossed --- a bound the record can confirm is there without fixing, from one look, just
where it lies. It is the first floor that is not a balance to be kept but a bound that cannot be passed --- notice only that even
a charge which balances runs up against a limit it does not get to cross, and that the instrument names the
limit as plainly as it names the tally. That candor will matter later. For now the record has run in a straight
line, the wheel's turning read as forward travel; but the wheel's own truth is that it closes on itself, and
the motion is about to bend back and shut into a loop. A loop keeps something a straight run cannot --- a
residue the closing leaves behind, which the next chapter is the first to read. The charge is entered and made
to balance; the record that must close is still open. The court does not sit until it does.

% <<<<<<<<<<<<<<<<<<<<<< APPEND ENDS <<<<<<<<<<<<<<<<<<<<<<
